% Generated by support/export-artifacts.py; do not edit.
% Slug: ch36-lanes
% Source: chapters/ch36.tex:32
% TikZ source: chapters/ch36.tex:34
% Original label: \label{fig:ch36-lanes}
% Original caption (LaTeX source):
%   Scalar code steps through elements one at a time while vector code covers four lanes per step, which pays only when lanes are independent.
% Caption: Scalar code steps through elements one at a time while vector code covers four lanes per step, which pays only when lanes are independent.
% Setup: support/compiler-pseudocode.sty (TikZ libraries and shared styles).
\begin{tikzpicture}[node distance=0.9cm and 1.2cm]
\node[basic block] (s1) {scalar\\ one element};
\node[basic block, below=of s1] (s2) {scalar\\ next element};
\node[basic block, right=of s1, minimum width=2.6cm] (v1) {vector\\ lanes 0--3};
\node[basic block, below=of v1, minimum width=2.6cm] (v2) {vector\\ lanes 4--7};
\draw[cfg edge] (s1) -- (s2);
\draw[cfg edge] (v1) -- (v2);
\end{tikzpicture}

% Original figure LaTeX:
% \begin{figure}[htbp]
% \centering
% \begin{tikzpicture}[node distance=0.9cm and 1.2cm]
% \node[basic block] (s1) {scalar\\ one element};
% \node[basic block, below=of s1] (s2) {scalar\\ next element};
% \node[basic block, right=of s1, minimum width=2.6cm] (v1) {vector\\ lanes 0--3};
% \node[basic block, below=of v1, minimum width=2.6cm] (v2) {vector\\ lanes 4--7};
% \draw[cfg edge] (s1) -- (s2);
% \draw[cfg edge] (v1) -- (v2);
% \end{tikzpicture}
% \caption{Scalar code steps through elements one at a time while vector code covers four lanes per step, which pays only when lanes are independent.}
% \label{fig:ch36-lanes}
% \end{figure}
